\scanpage{76}\mspage{2}
On a un candidat naturel (pour $\Sigma_{\deleted{\unsure{J|I}}}\inserted{=\Sigma_J}$, $\deleted{D_J}\inserted{D=D_J}$, $t\in\Sigma^\ast$ donnés) comme droite $T$ munie de (5), savoir
\[
T_0\simeq O_t(1) \tag*{\deleted{(7)}}
\]
à cela près qu'il y a \uline{deux} isom. naturels
\[
T_0^{\otimes2}=O_t(2)\simeq T_{\Sigma,t}\simeq O_t(\unsure{2|r})(\underline{\unsure{\omega|w}}),\tag*{\deleted{(8)}}
\]
correspondant aux deux éléments $w$ de $\underline\omega$, savoir
\[
\deleted{\unsure{t}}\ \xi\longmapsto\xi\wedge w \qquad
T_0^{\otimes2}\xrightarrow{\sim}O_t(2)\simeq T_{\Sigma,t}\simeq\underline O_t(2)(\underline\omega) \tag*{(7)}
\]

On va lever cette ambiguité, en remplaçant $T_0=O_t(1)$ par une variante tordue. Le tordre par un torseur sous $\mu_2$ \inserted{$=\mu_2(k)$} ne nous avance pas, car $\deleted{\unsure{O|T}}T_0^{\otimes2}$ n'aura pas changé (à isom can. près) on va le tordre par un torseur (qu'on précisera) $Q_0$ sous $\mu_4=\mu_4(k)$ (qui peut donc s'interpréter comme l'ens des sommets d'un \uline{carré}, dont $\mu_4$ est le groupe des rotations, avec $\{i,-i\}=\mu_4^\ast$ comme \uline{paire} de gén. privilégiés \ldots).

On posera donc
\[
T_0=\underline O_t(1)\wedge_{\mu_4}Q_0 \tag*{(8)}
\]
d'où
\[
T_0^{\otimes2}\simeq\underline O_t(2)\wedge_{\mu_2}D_{Q_0} \tag*{(9)}
\]
où
\[
D_{Q_0}=Q_0/\{\pm1\} \tag*{(10)}
\]
est l'ens des deux diagonales du carré $Q_0$. Supposons donc qu'on se donne un carré \inserted{combinatoire (``sommets'')} $Q_0$, avec des isom. \deleted{\unsure{Aut|Art} \unsure{$D_{Q_0}$|$Q_0$} \unsure{$\Sigma$|$Z$|$\simeq$} \unsure{entre|noter} le groupe des}
\[
\left\{\begin{aligned}
\underline\omega\,Q_0\quad &\text{(ens des orientations de $Q_0$)} &&\simeq\mu_4^\ast=\{i,-i\}\\
D_{Q_0}\quad &\text{(ens des diagonales de $Q_0$)} &&\simeq\underline\omega
\end{aligned}\right. \tag*{(11)}
\]
% archive sheet number: 75
